\scanpage{17}
\mspage{15}
particulier, on peut considérer que le cristal avec sa bi-structure supplémentaire permet de relier entre eux les groupes formels en les diverses caractéristiques auxquels il donne naissance $\ldots$

1.16. \uline{Un retour en arrière}. La définition 1.1. et le sorite 1.3. sont un petit peu canulés. Au lieu de prendre dans 1.0. pour $C$ la catégorie des \inserted{$R$-}flèches $S\longrightarrow T$ qui $\ldots$, il faut prendre la catégorie des $R$-flèches $U\longrightarrow T$ qui $\ldots$, où $U$ est un ouvert induit non précisé de $S$. De plus, \deleted{\unsure{il}} la définition n'est guère raisonnable alors que si la catégorie fibrée envisagée $\underline{F}$ est un ``champ'' pour la topologie de Zariski sur $(\operatorname{Sch})_{/R}$, i.e. si on peut y recoller flèches et objets. Ceci est nécessaire en tous cas pour pouvoir dans 1.3. définir la notion d'image inverse, la définition que j'y ai donnée n'étant raisonnable que si $S,S'$ sont affines. Dans le cas général, il faut se localiser sur $S$ et $S'$ pour se ramener à cette situation. Autrement on bute sur des canulars idiots de nature globale, comme le fait que sans restriction sur $S'\to S$, la construction \deleted{indiquée} \inserted{envisagée} de somme amalgamée fait sortir de la catégorie des préschémas $\ldots$ Il est probable qu'il y aura bien d'autres canulars de détail dans ces notes, mais, je pense, sans conséquence !

Il est évidemment tentant de vouloir interpréter les cristaux de modules comme faisceaux de modules sur un certain site. C'est possible, en prenant le ``\uline{site cristallogène} \uline{de} $S$ \uline{sur} $R$'', qui est précisément le site dont la catégorie sous-jacente \inserted{$C$} est celle des $R$-morphismes $U\longrightarrow T$ ($U$ ouvert de $S$, $U\longrightarrow T$ immersion fermée surjective définie par Ideal nilpotent sur $T$), la topologie est celle de Zariski: on prend comme familles couvrantes ``de définition'' de $(U\longrightarrow T)$ les familles définies par des recouvrements ouverts $(T_i)$ de $T$, chaque $T_i$ muni de la structure induite, et définissant $(U_i\longrightarrow T_i)$ par $U_i=U\cap T_i$. Un faisceau d'ensembles sur ce site s'identifie à un système de faisceaux d'ensembles $F_{U\to T}$ sur les \inserted{objets} buts des objets de $C$, avec, pour toute flèche $(U\longrightarrow T)\longrightarrow(U'\longrightarrow T')$ de $C$, un homomorphisme de l'image inverse de $F_{U'\to T'}$ \deleted{sur} \deleted{dans} \inserted{sur} $T$, dans $F_{U\to T}$ (homomorphisme qui n'est pas nécessairement un isomorphisme !)\inserted{(et qui est un \unsure{isom} si $T\unsure{\longrightarrow|-}T'$ est une \unsure{imm. ouverte})}. En particulier, associ-
\margin{16}
